\section{One-dimensional characters in subgroups of controlled index}
\label{l:characters-section}

A coset cap also controls a Fourier projection obtained from a
one-dimensional subgroup character.  This gives a different route to
the full isotypic estimates: first find one nonzero character space,
then average it over conjugates. The construction here assigns boxes
to rows or columns. Appendix~\ref{l:transfer-supplement} gives an
independent construction by successively removing longest lines,
together with its reciprocal-degree proof and transfer.

\begin{lemma}[A character from a row--column assignment]
\label{l:small-index-character}
For every irreducible complex representation of $S_n$, of degree $D$,
there are a subgroup $J\le S_n$ and a character
$\chi:J\to\{1,-1\}$ occurring in its restriction such that
\[
 [S_n:J]\le D^{1024}.
\]
\end{lemma}
\begin{proof}
We use the polytabloid realization of the Specht module of shape
$\lambda$; see \cite{sagan,etingof}.  Fix a labeling of its boxes. A row tabloid remembers which labels lie in each row while forgetting their order within that row; its polytabloid is the signed sum of its translates under the column stabilizer.
Assign every box either to its row or to its column.  Let $R$ permute
the assigned set in each row and let $C$ permute the assigned set in
each column.  These sets are disjoint, so $J=R\times C$.  Give $R$
the trivial character and each factor of $C$ the sign character.

Let $e$ be the polytabloid of the labeled diagram.  It transforms by
sign under the full column stabilizer, hence under $C$.  The average
$v=|R|^{-1}\sum_{r\in R}re$ is fixed by $R$ and still transforms by
sign under $C$.  It is nonzero.  Indeed the original row tabloid has
coefficient one in $e$, since the full row and column stabilizers
intersect trivially.  Every element of $R$ fixes that tabloid, so its
coefficient is also one in $v$.  This proves character occurrence for
every assignment.

It remains to choose the assignment with controlled index.  The arm
and leg of a box count boxes strictly to its right and strictly below
it.  Assign a box to its row when its arm is at least its leg, and to
its column otherwise.  Its hook length is at most twice the length
$l$ of the line receiving it.  If that line receives $a$ boxes, then
$a!\ge(a/e)^a$; we give a line with $a=0$ zero contribution below.
Consequently
\[
 \log\frac{\prod_{x\in\lambda}h(x)}
                  {\prod_{\text{assigned lines}}a!}
 \le n\log2+\sum_{\text{lines}}\{a+a\log(l/a)\}
 \le(\log2+1+2/e)n\le3n.
\]
Here $a\log(l/a)\le l/e$ and the total length of all rows and
columns is $2n$.  The hook-length formula therefore gives an assignment
with index
\begin{equation}\label{l:assignment-rough-index}
 I=[S_n:J]\le D e^{3n}.
\end{equation}

We convert the exponential factor into a power of $D$.  Put
$L=\max(\lambda_1,\lambda'_1)$.  Suppose first that $n\ge64$ and
$L\le3n/4$.  If $L\le n/8$, all hooks are at most $2L$, and
$D\ge(n/(2eL))^n\ge(4/e)^n\ge e^{n/64}$.
Otherwise transpose so that the first row has length $L$, and retain
a Young order ideal of $s=\lfloor L/4\rfloor$ boxes below it.
Such an ideal exists because $n-L\ge L/3$; moreover $s\ge n/64$.
For the retained diagram, let $c_j$ count its lower boxes in column
$j$.  Dividing its hook formula by the lower diagram's formula gives
\[
 D\ge\binom{L+s}s
       \prod_{j=1}^L\left(1+\frac{c_j}{L-j+1}\right)^{-1}
 \ge\binom{L+s}s\exp\left(-\frac{s}{L-s+1}\right).
\]
The first inequality also uses extension of tableaux from a subdiagram.
For the second, $c_j=0$ for $j>s$ and $\sum_jc_j=s$.
Since $s\le L/4$, the logarithm of the last expression is at least
$s\log5-s/(L-s+1)\ge s$.  Thus $\log D\ge n/64$ throughout
this case.  Equation~\eqref{l:assignment-rough-index} now gives
$I\le D^{193}$.

Suppose instead that $L>3n/4$.  Orient a longest line as the top row,
of length $L$, and write $\zeta$ for the lower diagram of size
$s=n-L$.  Assign the entire top row to its row and use the preceding
arm--leg assignment on $\zeta$.  Hook comparison and
\eqref{l:assignment-rough-index} applied to the tail give
\[
 I\le\binom nsD_\zeta e^{3s},\qquad
 D\ge\binom nsD_\zeta
              \exp\left(-\frac{s}{L-s+1}\right).
\]
If $s\ge1$, then $\binom ns\ge(n/s)^s\ge4^s$ and
$L-s+1\ge2$.  Hence
\[
 \log D\ge s(\log4-1/2),\qquad
 \log(I/D)\le\tfrac72s.
\]
These bounds imply $I\le D^{1024}$.  Transposing the assignment
back preserves the subgroup index and interchanges trivial and sign
factors.  If $s=0$, take the full group and its trivial or sign
character.  Finally, for $n<64$ and $D>1$, the trivial subgroup has
index $n!\le64^{64}<2^{1024}\le D^{1024}$.  Degree-one
representations again use the full group.  This covers every case.
\end{proof}

\subsection{From a character space to an isotypic space}

The index bound is useful because a subgroup character gives an
orthogonal projection through a signed average.  Positivity of the
original mass function controls that signed average.  The next lemma
also explains why repeated copies of a subgroup representation cause
no loss.

\begin{lemma}[Averaging character projections]
\label{l:subcharacter-isotypic}
Let $G$ be a finite group, $H\le G$, and let $f\ge0$ have mass
$z\le1$ and satisfy
$f(gH)\le B/[G:H]$ for every $g\in G$.  Let $\tau\in\Irr(H)$
have degree $d_\tau$.  Suppose $J\le H$ has a one-dimensional
unitary character $\chi$ occurring in $\tau|_J$, and put $I=[H:J]$.
For $\rho\in\Irr(G)$ of degree $D$, let $Q_\tau$ be the full
$\tau$-isotypic projection in $\rho|_H$.  Then
\begin{equation}\label{l:subcharacter-isotypic-bound}
 D\|\widehat f(\rho)Q_\tau\|_{\HS}^2
 \le BzI d_\tau.
\end{equation}
All copies of $\tau$ are included in $Q_\tau$.
\end{lemma}
\begin{proof}
Define the complex mass function $e_{J,\chi}$ to be
$\overline{\chi(j)}/|J|$ on $J$ and zero elsewhere.  Its Fourier
transform is the orthogonal projection $P$ onto vectors satisfying
$\rho(j)v=\chi(j)v$ for all $j\in J$.  With $U_H,U_J$ denoting
uniform probability masses on the subgroups, positivity gives
\[
 |f*e_{J,\chi}|\le f*U_J\le I(f*U_H)\le\frac{IB}{|G|},
 \qquad \sum_g|(f*e_{J,\chi})(g)|\le z.
\]
Plancherel therefore implies
$D\|\widehat f(\rho)P\|_{\HS}^2\le IBz$.
Every conjugate of $(J,\chi)$ inside $H$ satisfies the same bound.

Average the corresponding projections, writing
$\overline P=|H|^{-1}\sum_{h\in H}\rho(h)P\rho(h)^*$.
On the full $H$-isotypic space $V_\eta\otimes\mathcal A_\eta$,
a group-algebra element acts on $V_\eta$ tensored with the identity
on the multiplicity space $\mathcal A_\eta$.  Schur's lemma thus
gives
\[
 \overline P|_{V_\eta\otimes\mathcal A_\eta}
   =\frac{\operatorname{rank}(P|_{V_\eta})}{d_\eta}I.
\]
These scalars are nonnegative and the scalar for $\eta=\tau$ is
at least $1/d_\tau$.  Hence $\overline P\ge d_\tau^{-1}Q_\tau$.
Test this positive-operator inequality against
$\widehat f(\rho)^*\widehat f(\rho)$, and average the preceding
Plancherel bounds.  The result is
\eqref{l:subcharacter-isotypic-bound}.
\end{proof}

\begin{theorem}[Character-projection transfer]
\label{l:character-lift}
Let $H_i=\Sym(M_i)\le S_n$, $1\le i\le b$, for disjoint nonempty
sets $M_i$.  Let $f\ge0$ have mass $z\le1$ and suppose
\[
 f(gH_i)\le\frac B{[S_n:H_i]}
       \quad(g\in S_n,\ 1\le i\le b).
\]
For every irreducible $\rho$ of degree $D$,
\begin{equation}\label{l:character-lift-bound}
 \|\widehat f(\rho)\|_{\HS}^2
 \le BzbC_1D^{-1+1026/b},
\end{equation}
where $C_1=\sup_{m\ge1}\sum_{\tau\in\Irr(S_m)}d_\tau^{-1}$.
More generally, if $a\ge0$ and every type of every $H_i$ has a one-dimensional
subcharacter of index at most $d_\tau^a$, the exponent $1026/b$
in \eqref{l:character-lift-bound} may be replaced by $(a+2)/b$.
\end{theorem}
\begin{proof}
For each type $\tau$ of $H_i$, Lemma~\ref{l:subcharacter-isotypic}
and the assumed subcharacter give
\[
 D\|\widehat f(\rho)Q_{i,\tau}\|_{\HS}^2
       \le Bz d_\tau^{a+1}.
\]
Restrict $\rho$ to the product of the $H_i$.  In any occurring
joint type, the product of the factor dimensions is at most $D$,
including when that type has multiplicity greater than one.
Some factor therefore has dimension at most $R=D^{1/b}$.  The sum
of the commuting projections $Q_{i,\tau}$ with $d_\tau\le R$
dominates the identity.  Taking its trace against
$\widehat f(\rho)^*\widehat f(\rho)$ gives
\[
 D\|\widehat f(\rho)\|_{\HS}^2
 \le Bz\sum_{i=1}^b\sum_{d_\tau\le R}d_\tau^{a+1}
 \le BzbC_1R^{a+2}.
\]
The last inequality uses
$d_\tau^{a+1}\le R^{a+2}d_\tau^{-1}$.
Lemma~\ref{l:small-index-character} supplies $a=1024$.
\end{proof}

For example, take $b=8192$ and $B=2$.  If $f_n\le\mu_n$ has
mass tending to one and $\widehat\mu_n(\sgn)=o(1)$, then
\eqref{l:character-lift-bound} and Proposition~\ref{l:amplification}
give $\|\mu_n^{*4}-U_{S_n}\|_{\TV}\to0$.
Indeed $\varepsilon=1026/8192<1/4$, and the nonlinear Plancherel
sum for four factors is bounded by
$(2bC_1)^4\sum_{D_\lambda>1}D_\lambda^{1-4(1-\varepsilon)}$,
which tends to zero even using only the inverse-first-degree sum.
This statement requires only the displayed caps, mass and sign input.
